% ============================================================
% MODELO AUXILIAR: SOLICITAÇÃO DE PRORROGAÇÃO DE PRAZO
% Arquivo: modelos/solicitacao-prorrogacao.tex
%
% Peça auxiliar em formato de DIEx para solicitar à autoridade
% instauradora a prorrogação do prazo da sindicância.
% ============================================================

\pdfbookmark[1]{Solicitação de Prorrogação de Prazo}{solicitacao-prorrogacao-\getProrrogacaoNumero}

\ValidarSolicitacaoProrrogacao

\clearpage
\thispagestyle{empty}

\cabecalhoInstitucional

\vspace{3em}

\noindent\textbf{DIEx nº \getProrrogacaoNumero\ -- \getProrrogacaoOrigem}\\
\textbf{EB:} \getNup

\vspace{1em}

\begin{flushright}
\SeCampoPreenchido{prorrogacaoLocalData}{\getProrrogacaoLocalData}
\SeCampoVazio{prorrogacaoLocalData}{\localData}
\end{flushright}

\vspace{1em}

\noindent\textbf{Do} \getSindicantePosto\ \getSindicanteNome

\noindent\textbf{Ao Sr.} \getProrrogacaoDestinatario

\vspace{0.5em}

\noindent\textbf{Assunto:} prorrogação de prazo

\noindent\textbf{Referência:} \getProrrogacaoReferencia

\vspace{1em}

\ConteudoDocumentoFlexivel{%
\begin{enumerate}
  \item Trata o presente expediente sobre a solicitação de prorrogação de prazo para o prosseguimento das atividades da sindicância instaurada nesta OM, conforme referência.

  \item Solicito-vos a prorrogação do prazo da sindicância da qual sou encarregado por \getProrrogacaoPrazo, a contar de \getProrrogacaoDataInicio, com a finalidade de dar continuidade aos trabalhos realizados na sindicância.

  \item A solicitação é apresentada em \getProrrogacaoDataSolicitacao, com \getProrrogacaoAntecedenciaHoras\ horas de antecedência em relação ao término do prazo vigente, previsto para \getProrrogacaoTerminoPrazoAtual.

  \item Até o momento, foram realizados os seguintes atos: \getProrrogacaoAtosRealizados. A sindicância encontra-se na seguinte fase: \getProrrogacaoFaseAndamento.

  \item A prorrogação mostra-se necessária em razão de \getProrrogacaoJustificativa.
\end{enumerate}
}

\vspace{3em}

\begin{center}
\textbf{\getSindicanteNome} -- \getSindicantePosto\\
Sindicante
\end{center}

\clearpage
